\secciongrande{Programación de scripts: Dónde y cómo colocar los scripts}

\resumen{
  Los scripts de javascript se insertan en un documento html utilizando la etiqueta \texttt{<script>}. Pueden ubicarse tanto en la sección \texttt{<head>} como en el cuerpo \texttt{<body>}. La etiqueta de apertura debe incluir el atributo \texttt{type="text/javascript"} para indicar al navegador el formato del código.
}

\definicion{Etiqueta script}{Etiqueta html contenedora que designa un bloque de código javascript, comenzando con \texttt{<script>} y terminando con \texttt{</script>}.}

\subtituloA{Ubicación en el documento html}

\lista{
  \textbf{Encabezado del script:} Colocación entre \texttt{<head>} y \texttt{</head>}. |
  \textbf{Cuerpo del script:} Colocación entre \texttt{<body>} y \texttt{</body>}.
}
